Để cấu trúc mô hình phân phối đa biến từ các dữ liệu thực tế, nghiên cứu áp dụng phương pháp Hàm suy diễn cho các biên (Inference Functions for Margins - IFM). Đây là một kỹ thuật ước lượng hợp lý cực đại (Maximum Likelihood Estimation - MLE) theo hai giai đoạn, giúp giảm thiểu độ phức tạp tính toán so với việc ước lượng đồng thời toàn bộ tham số.
\begin{itemize}
    \item Giai đoạn 1 (Ước lượng biên): Sử dụng hàm phân phối tích lũy thực nghiệm (ECDF) để chuyển đổi chuỗi tỷ suất sinh lợi gốc của từng tài sản thành các chuỗi quan sát giả đồng nhất trên không gian $U[0,1]$. Việc sử dụng ECDF giúp bảo toàn tuyệt đối các đặc tính bất đối xứng và phân phối đuôi dày của dữ liệu thực tế mà không cần áp đặt bất kỳ giả định tham số nào lên phân phối biên.
    \item Giai đoạn 2 (Ước lượng tham số Copula): Dựa trên các dữ liệu đã được chuẩn hóa ở miền $U[0,1]$, tiến hành tối đa hóa hàm logarit của hàm hợp lý để tìm ra bộ tham số tối ưu cho từng hàm Copula ứng viên (ví dụ: ma trận tương quan và bậc tự do đối với t-Student Copula).
\end{itemize}
Sau khi ước lượng, tiêu chuẩn Cramér-von Mises dựa trên Copula thực nghiệm sẽ được sử dụng để kiểm định độ khớp mô hình (Goodness-of-Fit), từ đó chọn ra cấu trúc phụ thuộc mô tả chính xác nhất hành vi dữ liệu.